\section{Introduction}\label{sec:intro}

A route planner's estimate is a plan, and a plan that is consistently optimistic is still useful if the operator knows \emph{where} it is optimistic. The public trip records of Delhivery, an Indian road-freight carrier\footnote{Kaggle dataset ``Delhivery logistics dataset'' (user \texttt{devarajv88}); licence terms are in \texttt{DATA\_LICENSE.md} of the repository. We could not re-open the dataset page on the day of writing, so it is cited by this footnote and not in the bibliography.}, carry for every leg both a planned time recorded from OSRM~\cite{luxen2011real} and the realised time, so that question can be put to operational data directly rather than to a simulation.

The starting observation is one-sided\ifjournal{} (Fig.~\ref{fig:premise})\fi. Across \fzLegsTotal{} legs, 98.3\% exceed the recorded estimate and the median leg takes \fzMedianGapRatioDpOne{} times it. Noise would centre on one, so the error is systematic. A systematic error is not automatically \emph{localised}, though: it could be spread evenly, in which case a single correction factor removes it. This paper asks where it concentrates, how confidently that can be said, and what a model can do about it, and it reports the answers at the strength the data support rather than the strength a first analysis gave.

That last clause matters because a first analysis of these data was more flattering than the one reported here. A Welch test with false-discovery control found \fzBottlenecks{} corridors significantly slower than the network; a test that respects the dependence between legs of one recurring service finds that the same procedure on meaningless corridor labels returns about \fzAuditPlaceboDiscoveriesDpOne{} ``bottlenecks'' (Section~\ref{sec:audit}). A residual model first beat its baseline by 6.5\%; against a baseline without an unfair cold-start fallback the margin is 4.6\% for the model that is served, and the validation cut did not favour the model at all (Section~\ref{sec:prediction}). We make five contributions, the first being the one we think transfers:
\begin{enumerate}
\item \textbf{A way to measure a corridor audit's own error.} A placebo (real cell structure, no corridor effect) and a power simulation (known shifts planted into 30 corridors) show that the standard Welch audit is wrong in about \fzPowerWelchBhFdpShiftHalfPct{}\% of its discoveries on this network, and that the dependence-aware replacement is correctly sized but finds only \fzPowerByPermShiftOnePct{}\% of corridors even when a 2.7-fold overrun is planted (Section~\ref{sec:audit}). A zero from such a test cannot be read as absence.
\item \textbf{A dependence-aware audit and the evidence that corridor effects nonetheless persist.} Legs of one recurring service are correlated (design effect about 2); a permutation test whose unit is the (corridor, service) cell confirms \fzAuditRobustHeadlineSlower{} of the \fzBottlenecks{} slow corridors, yet corridor means correlate at \fzStabilitySpearmanDpTwo{} between the two halves of the period.
\item \textbf{Two remedies that do not rescue the ranking.} Empirical-Bayes shrinkage, with a normal or a heavy-tailed prior, leaves top-20 lists nearly as unstable across support floors (overlap \fzEbTopTwentyJaccardEbDpTwo{} against \fzEbTopTwentyJaccardRawDpTwo{}), because the cause is who is eligible; and the standard planning-time and buffer indices are less stable than the mean they summarise.
\item \textbf{A residual model over a fair baseline, with intervals and a serving proof.} A hierarchical back-off bar replaces a baseline with an unfair cold-start fallback; the reported model beats it (difference \fzReportedVsBarDeltaDpTwo{} minutes; day-block 95\% interval [\fzReportedVsBarCiLowDpTwo{}, \fzReportedVsBarCiHighDpTwo{}]), on all \fzSlicesTotal{} slices and in every rolling fold, mostly on corridors without history; the deployed Spark model is not robust across folds. The reported model is exact through the serving reader and the running stream.
\item \textbf{Two evaluation pitfalls and what they teach}: an optimiser default that disabled the loss the model was chosen for, and a replay that scored early events with later history (Section~\ref{sec:pitfalls}).
\end{enumerate}
We do not claim to beat the planner: the OSRM estimate has no access to corridor history, and a model that does is not a fair comparison. OSRM is reported as context throughout. The batch analyses sit inside a streaming serving path and a measured agent layer, described \ifjournal in Sections~\ref{sec:stream} and~\ref{sec:agents}\else in the extended version\fi; both evaluation pitfalls arose there, and the agents' data are synthetic.

\ifjournal
\begin{figure}[t]
\centering
\includegraphics[width=0.88\columnwidth]{fig1_gap_ratio_distribution.pdf}
\caption{Distribution of realised over recorded planned time across \fzLegsTotal{} legs (source: \texttt{benchmarks/raw/w1\_*}). The mass sits well to the right of 1.0: 98.3\% of legs run over the estimate and the median ratio is \fzMedianGapRatioDpOne{}.}
\label{fig:premise}
\end{figure}
\fi
